On nomme les évènements suivants :

\begin{itemize}
	\item $S$ : «~le client choisit un sorbet~»
	\item $L$ : «~le client choisit une glace au lait~»
	\item $G$ : «~le client choisit une crème glacée~»
	\item $C$ : «~le client choisit un cornet~»
\end{itemize}

	Ainsi $\overline{C}$ : «~le client choisit un pot~»

\begin{enumerate}
	\item $P\left(S\right) = \dfrac{28}{200} = \dfrac{7}{50}$.

	\item $P\left(S\cap \overline{C}\right) = \dfrac{12}{200} = \dfrac{3}{50}$.

	\item $P_{C}\left(L\right) = \dfrac{64}{160} = \dfrac{2}{5}$.

	\item $P_{G}\left(\overline{C}\right) = \dfrac{24}{104} = \dfrac{3}{13}$.
\end{enumerate}

\needspace{10\baselineskip}
